\documentclass[12pt,a4paper]{report}

%% Vložení užitečných balíčků
%%% Hlavní soubor. Zde se definují základní parametry a~odkazuje se na ostatní části. %%%

%% Verze pro jednostranný tisk:
% Okraje: levý 40mm, pravý 25mm, horní a~dolní 25mm
% (ale pozor, LaTeX si sám přidává 1in)
\setlength\textwidth{145mm}
\setlength\textheight{247mm}
\setlength\oddsidemargin{15mm}
\setlength\evensidemargin{15mm}
\setlength\topmargin{0mm}
\setlength\headsep{0mm}
\setlength\headheight{0mm}
% \openright zařídí, aby následující text začínal na pravé straně knihy
\let\openright=\clearpage

%% Pokud tiskneme oboustranně:
% \documentclass[12pt,a4paper,twoside,openright]{report}
% \setlength\textwidth{145mm}
% \setlength\textheight{247mm}
% \setlength\oddsidemargin{14.2mm}
% \setlength\evensidemargin{0mm}
% \setlength\topmargin{0mm}
% \setlength\headsep{0mm}
% \setlength\headheight{0mm}
% \let\openright=\cleardoublepage

\usepackage[pdfa,colorlinks=false,urlcolor=black,bookmarksopen=true]{hyperref} 	% Záložky v~souboru

%% Vytváříme PDF/A-2u
\usepackage[a-2u]{pdfx}

%% Přepneme na českou sazbu a~fonty Latin Modern
\usepackage[czech]{babel}
\usepackage{lmodern}
\usepackage[T1]{fontenc}
\usepackage{textcomp}

%% Použité kódování znaků: obvykle latin2, cp1250 nebo utf8:
\usepackage[utf8]{inputenc}
\usepackage{csquotes}
\MakeOuterQuote{"}

%%% Další užitečné balíčky (jsou součástí běžných distribucí LaTeXu)
\usepackage{amsmath}        % rozšíření pro sazbu matematiky
\usepackage{amsfonts}       % matematické fonty
\usepackage{amsthm}         % sazba vět, definic apod.
\usepackage{amssymb}
\usepackage{bbding}         % balíček s nejrůznějšími symboly
			    % (čtverečky, hvězdičky, tužtičky, nůžtičky, ...)
\usepackage{bm}             % tučné symboly (příkaz \bm)
\usepackage{graphicx}       % vkládání obrázků
\usepackage{fancyvrb}       % vylepšené prostředí pro strojové písmo
\usepackage{indentfirst}    % zavede odsazení 1. odstavce kapitoly
\usepackage[numbers]{natbib}         % zajištuje možnost odkazovat na literaturu
			    % stylem AUTOR (ROK), resp. AUTOR [ČÍSLO]
\usepackage[nottoc]{tocbibind} % zajistí přidání seznamu literatury,
                            % obrázků a~tabulek do obsahu
\usepackage{icomma}         % inteligetní čárka v~matematickém módu
\usepackage{dcolumn}        % lepší zarovnání sloupců v~tabulkách
\usepackage{booktabs}       % lepší vodorovné linky v~tabulkách
\usepackage{paralist}       % lepší enumerate a~itemize
\usepackage{xcolor}         % barevná sazba
\usepackage{float}
\usepackage{ifthen}
\usepackage{cancel}
\usepackage{enumitem}
\usepackage{ragged2e}
\usepackage{needspace}		
\usepackage{wrapfig}		% plovoucí obrázky nalevo/napravo
\usepackage{pgfplots}
\usepackage{subcaption}
\pgfplotsset{compat=1.15}
\usepackage{mathrsfs}
\usetikzlibrary{arrows}

\justifying
\graphicspath{{./components/images/}}

%% Definice různých užitečných maker (viz popis uvnitř souboru)

% Toto makro definuje kapitolu, která není očíslovaná, ale je uvedena v~obsahu.
\def\chapwithtoc#1{
\chapter*{#1}
\addcontentsline{toc}{chapter}{#1}
}


\DefineBibliographyStrings{czech}{ % Definování vlastního názvu pro bibliografii v češtině
  bibliography = Seznam použité literatury,
}

% TODO
\newcommand{\todo}[1]{\textcolor{red}{(\noindent TODO: #1.)}}


% Celá jména
\newcommand{\name}[1]{\mbox{\textsc{#1}}}


% Cesty
\newcommand{\appendixpath}{components/appendix}

% Images
\newcommand{\subfigwidth}{6cm}        % šířka okénka pro "podobrázky"
\newcommand{\wrappedfigwidth}{6cm}
\newcommand{\fullhd}{0.2}             % měřítko Full HD obrázku
\newcommand{\normalipe}{0.8}          % standardní měřítko IPE obrázku
\newcommand{\fractalscale}{0.3}       % měřítko obrázku softwarem vygenerovaného fraktálu

%%% Prostředí pro sazbu kódu, případně vstupu/výstupu počítačových
%%% programů. (Vyžaduje balíček fancyvrb -- fancy verbatim.)

\DefineVerbatimEnvironment{code}{Verbatim}{fontsize=\small, frame=single}

% Zapne černé "slimáky" na koncích řádků, které přetekly, abychom si
% jich lépe všimli.
\overfullrule=1mm

% Trochu volnější nastavení dělení slov, než je default.
\lefthyphenmin=2
\righthyphenmin=2

\setlength{\parskip}{0.3em}



%%% Údaje o~práci

%%% Název práce v~jazyce práce (přesně podle zadání)
\def\NazevPrace{Virtualizační platforma Proxmox}

%%% Název práce v~angličtině
\def\NazevPraceEN{\todo{Title}}

%%% Typ práce
\def\TypPrace{Seminární práce}

%%% Jméno autora
\def\AutorPrace{Martin Kopecký}



%%% Rok odevzdání
\def\RokOdevzdani{2024}

%%% Název katedry nebo ústavu, kde byla práce oficiálně zadána
%%% (dle Organizační struktury MFF UK, případně plný název pracoviště mimo MFF)
% \def\Katedra{\todo{Katedra}}
% \def\KatedraEN{\todo{Department}}

% Jedná se o~katedru (department) nebo o~ústav (institute)?
% \def\TypPracoviste{Katedra}
% \def\TypPracovisteEN{Department}

% Vedoucí práce: Jméno a~příjmení s~tituly
% \def\Vedouci{\todo{Vedouci}}

% Pracoviště vedoucího (opět dle Organizační struktury MFF)
% \def\KatedraVedouciho{\todo{Katedra}}
% \def\KatedraVedoucihoEN{\todo{Department}}

% Studijní program a~obor
\def\StudijniProgram{Aplikovaná informatika}
\def\StudijniObor{Informační systémy}

%%% Nepovinné poděkování (vedoucímu práce, konzultantovi, tomu, kdo
%%% zapůjčil software, literaturu apod.)
%\def\Podekovani{%
%\todo{Doplnit poděkování}
%\input{acknowledgements.txt}
%}

%%% Abstrakt (doporučený rozsah cca 80-200 slov; nejedná se o~zadání práce)
%\def\Abstrakt{%
%\todo{Doplnit abstrakt (česky)}
%\input{abstract_cz.txt}
%}
%\def\AbstraktEN{%
%\todo{Doplnit abstrakt (anglicky)}
%\input{abstract_en.txt}
%}

%%% 3 až 5 klíčových slov (doporučeno), každé uzavřeno ve složených závorkách
\def\KlicovaSlova{%
{}, {klicova slova}
}
%\def\KlicovaSlovaEN{%
%{keywords}
%}

%%% Balíček hyperref, kterým jdou vyrábět klikací odkazy v~PDF,
%%% ale hlavně ho používáme k~uložení metadat do PDF (včetně obsahu).
%%% Většinu nastavítek přednastaví balíček pdfx.
\hypersetup{unicode}
\hypersetup{breaklinks=true}


\renewcommand{\bibname}{Seznam použité literatury}
\addbibresource{99-literature/literature.bib}

%% Titulní strana a~různé povinné informační strany
\begin{document}
%%% Titulní strana práce a další povinné informační strany

%%% Titulní strana práce

\pagestyle{empty}
\hypersetup{pageanchor=false}

\begin{center}

%\centerline{\mbox{\includegraphics[width=166mm]{components/images/ujep.png}}}
\null


{\bf\Large Univerzita Jana Evangelisty Purkyně \\ v Ústí nad Labem}

\vspace{2mm}

{\bf\Large Přírodovědecká fakulta}

\vspace{20mm}

\includegraphics[width=80mm]{ujep.png}

\vspace{20mm}


{\LARGE\bfseries\NazevPrace}

{\large\bfseries\TypPrace}

\vfill


\begin{flushleft}
\begin{tabular}{>{\bfseries}r l}
    \textbf{Vypracoval:} & \AutorPrace \\
    % Vedoucí práce:\Vedouci 
    % Konzultant: \Vedouci
    \\

    \textbf{Studijní program:}  &  \StudijniProgram \\
    \textbf{Studijní obor:}  &   \StudijniObor \\
\end{tabular}
\end{flushleft}

\vfill

% Zde doplňte rok
\large{Ú\raisebox{-0.85ex}{\textsuperscript{STÍ NAD}} L\raisebox{-0.85ex}{\textsuperscript{ABEM}} \RokOdevzdani}

\end{center}

%\newpage

%%% Následuje vevázaný list -- kopie podepsaného "Zadání diplomové práce".
%%% Toto zadání NENÍ součástí elektronické verze práce, nescanovat.

%%% Strana s čestným prohlášením k diplomové práci

%\openright
%\hypersetup{pageanchor=true}
%\pagestyle{plain}
%\pagenumbering{roman}
%\vglue 0pt plus 1fill

%\noindent
%Prohlašuji, že jsem tuto diplomovou práci vypracoval(a) samostatně a výhradně
%s~použitím citovaných pramenů, literatury a dalších odborných zdrojů.
%Tato práce nebyla využita k získání jiného nebo stejného titulu.
%
%\medskip\noindent
%Beru na~vědomí, že se na moji práci vztahují práva a povinnosti vyplývající
%ze zákona č. 121/2000 Sb., autorského zákona v~platném znění, zejména skutečnost,
%že Univerzita Karlova má právo na~uzavření licenční smlouvy o~užití této
%práce jako školního díla podle §60 odst. 1 autorského zákona.
%
%\vspace{10mm}
%
%\hbox{\hbox to 0.5\hsize{%
%V \hbox to 6em{\dotfill} dne \hbox to 6em{\dotfill}
%\hss}\hbox to 0.5\hsize{\dotfill\quad}}
%\smallskip
%\hbox{\hbox to 0.5\hsize{}\hbox to 0.5\hsize{\hfil Podpis autora\hfil}}
%
%\vspace{20mm}
%\newpage

%%% Poděkování

%\openright
%
%\noindent
%\Podekovani
%
%\newpage

%%% Povinná informační strana diplomové práce

%\openright
%
%\vbox to 0.5\vsize{
%\setlength\parindent{0mm}
%\setlength\parskip{5mm}
%
%Název práce:
%\NazevPrace
%
%Autor:
%\AutorPrace
%
%\TypPracoviste:
%\Katedra
%
%Vedoucí diplomové práce:
%\Vedouci, \KatedraVedouciho
%
%Abstrakt:
%\Abstrakt
%
%Klíčová slova:
%\KlicovaSlova
%
%\vss}\nobreak\vbox to 0.49\vsize{
%\setlength\parindent{0mm}
%\setlength\parskip{5mm}
%
%Title:
%\NazevPraceEN
%
%Author:
%\AutorPrace
%
%\TypPracovisteEN:
%\KatedraEN
%
%Supervisor:
%\Vedouci, \KatedraVedoucihoEN
%
%Abstract:
%\AbstraktEN
%
%Keywords:
%\KlicovaSlovaEN
%
%\vss}

\newpage

\openright
\pagestyle{plain}
\pagenumbering{arabic}
\setcounter{page}{1}

%%% Strana s automaticky generovaným obsahem bakalářské práce

\tableofcontents

%%% Jednotlivé kapitoly práce jsou pro přehlednost uloženy v~samostatných souborech
\chapwithtoc{Předmluva}\label{chapter:predmluva}
\todo{napsat předmluvu}

\cite{Prusinkiewicz1990}
aasda
\chapwithtoc{Historie}\label{chapter:predmluva}

Proxmox je oblíbená virtualizační platforma, která kombinuje virtuální servery (Proxmox Virtual Environment, PVE) a systém pro správu clusterů a virtualizovaného úložiště (Proxmox Virtual Environment Cluster). Její historie sahá do roku 2008, kdy byla poprvé uvedena na trh společností Proxmox Server Solutions GmbH, založenou v Rakousku.

Proxmox VE byl vytvořen jako open-source řešení, které poskytuje efektivní a snadno použitelné prostředí pro správu virtuálních serverů. Od začátku byl postaven na Linuxovém jádru a využíval virtualizační technologie jako KVM (Kernel-based Virtual Machine) a kontejnerovou technologii LXC (Linux Containers), což umožňovalo uživatelům vytvářet a spravovat virtuální stroje a kontejnery z jediného rozhraní.

V průběhu let došlo k několika důležitým vývojovým krokům. Významným milníkem bylo zavedení Proxmox VE Cluster, což umožnilo správu více serverů jako jednoho celku, což výrazně zjednodušilo správu zdrojů a redundanci. Další významnou aktualizací bylo zavedení Ceph, rozšířeného open-source úložného řešení, které umožňuje vysoce škálovatelné blokové úložiště.

Proxmox se také vyznačuje svým uživatelským rozhraním založeným na webovém prohlížeči, které umožňuje jednoduché správu a monitoring virtuálních strojů a úložiště. Uživatelé oceňují jeho flexibilitu a širokou škálu konfiguračních možností, které jsou přístupné i bez hlubokých znalostí Linuxu.

Proxmox VE je také známý svou komunitou a širokou podporou. Jako open-source řešení má aktivní komunitu vývojářů a uživatelů, kteří neustále přispívají k jeho vývoji a zlepšování. To zahrnuje pravidelné aktualizace, bezpečnostní opravy a nové funkce, což udržuje Proxmox VE jako relevantní a konkurenceschopnou platformu v oblasti virtualizace.

V současnosti je Proxmox široce používán v mnoha odvětvích, od malých podniků až po velké korporace, díky své schopnosti efektivně spravovat rozsáhlé a komplexní virtualizované prostředí. Jeho kombinace výkonu, flexibility a nízkých nákladů činí z Proxmoxu atraktivní volbu pro mnoho organizací hledajících efektivní virtualizační řešení.
% \include{\chapterpath{01}/ch01_limity.tex}

%%% Seznam použité literatury
\printbibliography

%%% Obrázky v~bakalářské práci
%%% (pokud jich je malé množství, obvykle není třeba seznam uvádět)
%\listoffigures

%%% Tabulky v~bakalářské práci (opět nemusí být nutné uvádět)
%%% U matematických prací může být lepší přemístit seznam tabulek na začátek práce.
% \listoftables

%%% Přílohy k~bakalářské práci, existují-li. Každá příloha musí být alespoň jednou
%%% odkazována z vlastního textu práce. Přílohy se číslují.
%%%
%%% Do tištěné verze se spíše hodí přílohy, které lze číst a~prohlížet (dodatečné
%%% tabulky a~grafy, různé textové doplňky, ukázky výstupů z počítačových programů,
%%% apod.). Do elektronické verze se hodí přílohy, které budou spíše používány
%%% v~elektronické podobě než čteny (zdrojové kódy programů, datové soubory,
%%% interaktivní grafy apod.). Elektronické přílohy se nahrávají do SISu a~lze
%%% je také do práce vložit na CD/DVD. Povolené formáty souborů specifikuje
%%% opatření rektora č. 72/2017.

\appendix

\openright
\end{document}